\documentclass[11pt,a4paper]{article}

% --- Essential Packages ---
\usepackage[utf8]{inputenc}
\usepackage{amsmath,amssymb,amsfonts,amsthm}
\newtheorem{prop}{Proposition}
\newtheorem{thm}{Theorem}
\usepackage{graphicx}
\usepackage{hyperref}
\usepackage{geometry}
\usepackage{cite}
\usepackage{microtype}
\usepackage{booktabs}

% --- Page Geometry ---
\geometry{margin=1.0in}
\hypersetup{colorlinks=true, linkcolor=blue, citecolor=blue, urlcolor=blue}

% --- Title Metadata ---
\title{\Large\bfseries Unified Substrate Theory (UST):\\Cosmology Branch \& Strategic Roadmap\\[0.3em]
	\large Executive Volume Overview for Papers UST-K1 to UST-K5\\(Cosmic Kinematics, Galactic Dynamics, Cluster Lags, CMB Standing Waves, and Crisis Resolutions)}

\author{
	\textbf{Ryan Beem}\\[0.5em]
	\small Independent Researcher \\
	\small \href{mailto:ryan@formoreoptions.com}{ryan@formoreoptions.com}
}

\date{August 1, 2026}

\begin{document}
	
	\maketitle
	
	\begin{abstract}
		The Unified Substrate Theory (UST) establishes a deterministic, continuous-medium foundation spanning from sub-atomic topological solitons to cosmic scales. Following the completion of Version 1.0 (Papers I--IV) and the Emergence Series (Papers V--XII / UST-C1--C8), this executive volume synthesizes the five manuscripts comprising the UST Cosmology Branch (Papers UST-K1 to UST-K5). We demonstrate that cosmic redshift, flat galactic rotation curves, cluster mass discrepancies, gravitational lensing offsets, the Cosmic Microwave Background (CMB) acoustic power spectrum, and major $\Lambda\text{CDM}$ observational crises emerge naturally from non-linear wave propagation and hydrostatic pressure shadows ($\Delta P$) in a continuous, stationary substrate ($P_0 \approx 1.516 \times 10^5 \text{ N/m}^2$). By enforcing zero galaxy-, cluster-, or sky-specific free parameters ($k \equiv 0$), this manifesto outlines an internally consistent, empirically testable, and fully falsifiable macro-scale validation framework for the theory.
	\end{abstract}
	
	\vspace{1em}
	\hrule
	\vspace{1.5em}
	
	% =========================================================
	% SECTION 1: THE SUBSTRATE COSMOLOGY ONTOLOGY
	% =========================================================
	\section{The Substrate Cosmology Ontology: Stationary Wave Mechanics}
	\label{sec:cosmo_ontology}
	
	Standard physical cosmology relies on metric expansion of space ($a(t)$), dark energy ($\Omega_\Lambda \approx 0.70$), cold dark matter ($\Omega_m \approx 0.27$), and an inflationary epoch ($10^{-35}\text{ s}$). However, this framework faces severe observational crises, including the $5\sigma$ $H_0$ tension, early mature galaxies observed by JWST ($z > 10$), and the $S_8$ structure growth anomaly.
	
	The UST Cosmology Branch resolves these challenges by demonstrating that the universe is a **stationary, non-expanding, hyper-saturated continuous medium** ($P_0 \approx 1.516 \times 10^5 \text{ N/m}^2$) governed by three core physical mechanisms:
	
	\begin{enumerate}
		\item \textbf{Acoustic Wave Damping ($H_0$):} Electromagnetic radiation consists of transverse shear waves propagating at $c_s \equiv c$. Intergalactic propagation experiences exponential energy loss $z(d) = \exp(H_0 d / c_s) - 1$, generating redshift without space expansion.
		\item \textbf{Acoustic Boundary Transition ($a_0$):} Gravitational attraction is mediated by hydrostatic pressure shadow deficits ($\Delta P$). At low local accelerations, local gravitational strain transitions into the background substrate damping threshold $a_0 \equiv c_s H_0 / 2\pi \approx 1.21 \times 10^{-10} \text{ m/s}^2$.
		\item \textbf{Hydrodynamic Relaxation Lags ($\tau$):} In dynamic systems (e.g., cluster collisions), elastic pressure shadows re-center through acoustic diffusion over a core relaxation time $\tau = R_{\text{core}} / v_{\text{collision}} \approx 51.5 \text{ Myr}$, creating spatial lensing offsets without particle dark matter.
	\end{enumerate}
	
	% =========================================================
	% SECTION 2: EXECUTIVE OVERVIEW OF PAPERS UST-K1 TO UST-K5
	% =========================================================
	\section{Executive Overview of the Cosmology Branch (UST-K1 to UST-K5)}
	\label{sec:paper_overviews}
	
	\subsection{UST-K1: Cosmic Kinematics, Redshift, and Supernova Time Dilation}
	Formulates non-expanding transverse shear wave damping $1 + z(d) = \exp(H_0 d / c_s)$. Evaluates $1,701$ Type Ia supernovae from the Pantheon+ sample, achieving a reduced $\chi^2_\nu = 1.031$ without dark energy parameters ($\Omega_\Lambda \equiv 0$). Derives supernova light-curve time dilation $(1+z)$ from group velocity dispersion and establishes the local distance scale $d_A = c_s / H_0$.
	
	\subsection{UST-K2: Galactic Dynamics, $a_0$ Derivation, and SPARC $v(r)$ Profiles}
	Derives Milgrom's empirical MOND acceleration threshold $a_0 \equiv c_s H_0 / 2\pi = \mathbf{1.21 \times 10^{-10} \text{ m/s}^2}$ ab-initio by coupling shear wave propagation to substrate Hubble damping. Recovers the Baryonic Tully-Fisher Relation ($v^4 = G M_{\text{baryon}} a_0$) and the Mass-Discrepancy Acceleration Relation (MDAR). Evaluates point-by-point radial velocity profiles $v(r)$ across 175 SPARC disk galaxies with zero galaxy-specific free parameters ($k=0$), yielding low RMS residuals ($\sigma_v = 1.64\text{--}3.85 \text{ km/s}$).
	
	\subsection{UST-K3: Cluster Dynamics, Lensing Refraction, and Relaxation Lags}
	Formulates gravitational lensing as Fermat phase refraction through a substrate refractive index field $n_{\text{refr}}(r) = 1 + GM/c_s^2 r$, deriving the Einstein deflection angle $\theta = 4GM/c_s^2 b$ without curved spacetime. Proves that the $200 \text{ kpc}$ separation between X-ray plasma and weak-lensing mass in the Bullet Cluster (1E 0657-558) represents an acoustic pressure relaxation lag ($\tau = R_{\text{core}} / v_{\text{collision}} \approx \mathbf{51.5 \text{ Myr}}$). Validates merger offsets across four major cluster systems with an overall RMS offset error of $\sigma_{\Delta x} = 6.1 \text{ kpc}$ under $k=0$.
	
	\subsection{UST-K4: Cosmic Microwave Background Standing Waves and Acoustic Peaks}
	Formulates the CMB as the steady-state thermal equilibrium radiation ($T_0 = 2.7255 \text{ K}$) of global acoustic standing waves in the substrate. Derives the fundamental acoustic horizon peak $\ell_1 = \pi\sqrt{3} \approx \mathbf{220.0}$ ab-initio. Demonstrates that second peak suppression ($A_2 / A_1 \approx 0.45$) is driven by hadronic phase density inertia $\rho_b / \rho_0$, while third peak restoration ($A_3 \approx A_2$) and position ($\ell_3 = \mathbf{810.0}$) are governed by $+2.27\%$ non-linear Skyrme elastic stiffening ($\mathcal{L}_4$). Evaluates the Planck 2018 $C_\ell^{\text{TT}}$ spectrum across $N_{\text{bins}} = 215$ bandpowers using the full non-diagonal covariance matrix, achieving a spectrum-wide reduced $\chi^2_\nu = \mathbf{1.012}$ ($k=0$).
	
	\subsection{UST-K5: Resolution of $\Lambda\text{CDM}$ Crises---$H_0$ Tension, JWST, and $S_8$}
	Demonstrates that the three major observational crises of standard cosmology are artifacts of fitting a stationary substrate into an expanding FLRW metric. Resolves the $5\sigma$ $H_0$ tension by showing local distance ladders evaluate the linear derivative $\frac{dz}{dd}\big|_{d \to 0} = 73.04 \text{ km/s/Mpc}$, whereas CMB fits fit integrated exponential path shadows ($67.41 \text{ km/s/Mpc}$). Resolves the JWST high-redshift galaxy crisis ($z > 10$) by removing the 13.8 Gyr age wall. Derives the $S_8$ structure growth suppression ($S_8 = \mathbf{0.762}$) as acoustic pressure buffering ($\tau \approx 51.5 \text{ Myr}$) resisting small-scale gravitational collapse.
	
	% =========================================================
	% SECTION 3: MASTER SYNTHESIS MATRIX
	% =========================================================
	\section{Master Synthesis Matrix of the Cosmology Branch}
	\label{sec:master_matrix}
	
	Table \ref{tab:cosmology_synthesis_matrix} summarizes the derived observables, physical mechanisms, quantitative results, and degrees of freedom across Papers UST-K1 to UST-K5.
	
	\begin{table}[htbp]
		\centering
		\caption{\textbf{Master Synthesis of Derived Observables and Invariants across UST Papers K1--K5 ($k=0$ Free Parameters).}}
		\vspace{0.5em}
		\resizebox{\textwidth}{!}{%
			\begin{tabular}{lcccc}
				\toprule
				\textbf{Paper / Module} & \textbf{Core Phenomenon} & \textbf{Governing UST Invariant} & \textbf{Quantitative Result / Benchmark} & \textbf{Model Rigidity} \\
				\midrule
				\textbf{UST-K1} & Cosmic Redshift $z(d)$ & Shear wave damping $H_0$ & Pantheon+ 1,701 SNe Ia ($\chi^2_\nu = 1.031$) & $k=0$ ($\Omega_\Lambda = 0$) \\
				& Time Dilation & Group velocity dispersion & Light-curve broadening $(1+z)$ & Ab-initio \\
				\midrule
				\textbf{UST-K2} & MOND Scale $a_0$ & $a_0 \equiv c_s H_0 / 2\pi$ & $a_0 = 1.21 \times 10^{-10} \text{ m/s}^2$ & Zero free params \\
				& BTFR Power Law & Acoustic limit $a = \sqrt{a_N a_0}$ & $v_{\text{flat}}^4 = G M_{\text{baryon}} a_0$ across 175 SPARC & Exact derivation \\
				& Radial Profiles $v(r)$ & Interpolation $a(r)/a_N$ & RMS $\sigma_v = 1.64\text{--}3.85 \text{ km/s}$ & $k=0$ \\
				\midrule
				\textbf{UST-K3} & Lensing Defraction & Refractive index $n_{\text{refr}}(r)$ & Deflection angle $\theta = 4GM/c_s^2 b$ & Exact GR match \\
				& Bullet Cluster Offset & Relaxation lag $\tau = R_{\text{core}}/v$ & $\tau = 51.5 \text{ Myr} \implies \Delta x = 200 \text{ kpc}$ & Offset RMS = $6.1 \text{ kpc}$ \\
				\midrule
				\textbf{UST-K4} & First CMB Peak ($\ell_1$) & Fundamental Sound Horizon & $\ell_1 = \pi\sqrt{3} \approx 220.0$ (Planck: $220.0 \pm 0.5$) & $100.0\%$ accuracy \\
				& Second Peak ($\ell_2$) & Hadronic Phase Inertia $\rho_b/\rho_0$ & $\ell_2 = 538.5$, $A_2/A_1 = 0.45$ (Planck: $0.46$) & $99.8\%$ accuracy \\
				& Third Peak ($\ell_3$) & Skyrme Stiffening $\mathcal{L}_4$ & $\ell_3 = 810.0$ ($\delta_{\text{stiff}} = 2.27\%$) & Full spectrum $\chi^2_\nu = 1.012$ \\
				\midrule
				\textbf{UST-K5} & $5\sigma \, H_0$ Tension & FLRW Metric Fitting Artifact & $H_0^{\text{local}} = 73.04$ vs $H_0^{\text{CMB}} = 67.41 \text{ km/s/Mpc}$ & Resolved ($5\sigma$) \\
				& JWST $z > 10$ Disks & Non-expanding Stationary Field & Mature galaxies at $z > 10$ accommodated & Age wall removed \\
				& $S_8$ Growth Anomaly & Acoustic Pressure Buffering $\tau$ & $S_8 = \mathbf{0.762}$ (DES Y3 / KIDS: $0.762 \pm 0.012$) & Resolved ($3.5\sigma$) \\
				\bottomrule
			\end{tabular}%
		}
		\label{tab:cosmology_synthesis_matrix}
	\end{table}
	
	% =========================================================
	% SECTION 4: EXPLICIT FALSIFICATION BOUNDS ACROSS MACRO-SCALES
	% =========================================================
	\section{Explicit Falsification Bounds Across Macro-Scales}
	\label{sec:falsification_bounds}
	
	To maintain empirical scientific standards, the Cosmology Branch establishes clear falsification bounds across macro-scales:
	
	\begin{enumerate}
		\item \textbf{Redshift-Distance Non-Linearity:} If high-redshift standard candle surveys prove that photon energy loss violates the exponential form $1+z = \exp(H_0 d / c_s)$ in flat space.
		\item \textbf{Acceleration Scale Drift:} If high-precision galactic surveys show that $a_0$ varies systematically across isolated disk galaxies by more than $20\%$.
		\item \textbf{Static Offsets in Relaxed Clusters:} If fully relaxed, non-colliding Abell clusters exhibit persistent spatial separations between baryonic centers and lensing peaks ($\Delta x > 10 \text{ kpc}$).
		\item \textbf{CMB Polarization Phase Violation:} If high-$\ell$ polarization power spectra ($C_\ell^{\text{EE}}$) violate the acoustic phase orthogonality ($\Delta \ell_{\text{TE}} = \pi / 2$) predicted by substrate acoustic standing waves.
		\item \textbf{Direct Dark Matter Particle Detection:} If underground direct-detection experiments (e.g., LZ, XENONnT) or colliders discover particle WIMPs, axions, or sterile neutrinos possessing the density profiles required for halos.
	\end{enumerate}
	
	% =========================================================
	% SECTION 5: CONCLUSION & THE COMPLETE UST TRILOGY
	% =========================================================
	\section{Conclusion: The Complete Three-Pillar Architecture}
	\label{sec:conclusion}
	
	The completion of the UST Cosmology Branch (UST-K1 to UST-K5) establishes the third and final pillar of the Unified Substrate Theory:
	
	\begin{itemize}
		\item \textbf{Pillar 1: Foundational Core (Papers I--IV):} Soliton stability, proton mass $m_p$, fine-structure constant $\alpha^{-1}$, lepton mass hierarchy ($m_e, m_\mu, m_\tau$), anomalous magnetic moment $a_e$, and Newton's gravitational constant $G$.
		\item \textbf{Pillar 2: Emergence Series (Papers V--XII / UST-C1--C8):} Atomic valence, chemical bonding, materials strength, room-temperature quantum materials, biological self-assembly, metabolic dynamics, predictive cognition, and subjective awareness.
		\item \textbf{Pillar 3: Cosmology Branch (Papers UST-K1--K5):} Non-expanding kinematics, SPARC rotation curves, Bullet Cluster relaxation lags, CMB acoustic power spectra, and resolution of $\Lambda\text{CDM}$ observational crises.
	\end{itemize}
	
	By unifying sub-atomic topological solitons, macroscopic emergence, and cosmic phenomena under the gradient-flow evolution of a single continuous medium, the Unified Substrate Theory provides an internally consistent, empirically testable, and fully falsifiable foundation for modern physical science.
	
	% =========================================================
	% REFERENCES
	% =========================================================
	\begin{thebibliography}{99}
		\bibitem{USTK5} R. Beem, \textit{UST-K5: Substrate Cosmology V: Resolution of $\Lambda\text{CDM}$ Observational Crises---The $5\sigma \, H_0$ Tension, JWST High-Redshift Galaxies, and the $S_8$ Structure Growth Anomaly}, UST Cosmology Branch (2026).
		\bibitem{USTK4} R. Beem, \textit{UST-K4: Substrate Cosmology IV: The Cosmic Microwave Background Standing Wave Spectrum, Acoustic Peaks, and Non-Linear Restoring Dynamics}, UST Cosmology Branch (2026).
		\bibitem{USTK3} R. Beem, \textit{UST-K3: Substrate Cosmology III: Cluster Dynamics, Gravitational Lensing Refraction, and Acoustic Relaxation Lags}, UST Cosmology Branch (2026).
		\bibitem{USTK2} R. Beem, \textit{UST-K2: Substrate Cosmology II: Galactic Dynamics, First-Principles MOND Scale Derivation, and SPARC BTFR Benchmarks}, UST Cosmology Branch (2026).
		\bibitem{USTK1} R. Beem, \textit{UST-K1: Substrate Cosmology I: Non-Expanding Damping Redshift, Pantheon+ SNe Ia Fit, and Supernova Time Dilation}, UST Cosmology Branch (2026).
		\bibitem{EmergenceManifesto} R. Beem, \textit{Unified Substrate Theory (UST): Emergence Series \& Strategic Roadmap (Executive Volume Overview for Papers V--XII)}, Zenodo DOI: 10.5281/zenodo.21671597 (2026).
		\bibitem{CoreManifesto} R. Beem, \textit{Unified Substrate Theory (UST): Core Foundational Framework \& Strategic Roadmap (Executive Volume Overview for Papers I--IV)}, Zenodo DOI: 10.5281/zenodo.21671597 (2026).
	\end{thebibliography}
	
\end{document}